% !TEX root = ../notes.tex

\documentclass[../notes.tex]{subfiles}

\begin{document}

\section{May 4}
Here we go.

\subsection{Archimedean Whittaker Functionals}
Let's retell our story of Whittaker models in the archimedean setting.
\begin{definition}[Whittaker functional]
	Fix a split reductive group $G$ over $\RR$, and let $(V,\pi)$ be a unitary representation. Given a character $\psi\colon\RR\to S^1$, we define $\psi_N$ on the unipotent radical of a Borel subgroup as before, by summing over simple root spaces. Then a \textit{Whittaker functional} is a linear map $\Lambda\colon V\to\CC$ for which
	\[\Lambda(\pi(u)v)=\psi_N(u)\Lambda(v).\]
	We will also ask that $\Lambda$ has moderate growth, in the sense of \Cref{def:mod-growth} below.
\end{definition}
It turns out that uniqueness is a subtle issue here.
\begin{exe}
	We compute the Whittaker functionals (without moderate growth) of the discrete series representation $D_k^+$ of the group $\op{SL}_2(\RR)$.
\end{exe}
\begin{proof}
	Let $V$ be the discrete series representation $D_k^+$ of $\op{SL}_2(\RR)$. By Frobenius reciprocity, a choice of Whittaker functional produces a map $D_k^+\into\op{Mor}_{\mathrm{Set}}((N(\RR),\psi_N)\backslash\op{SL}_2(\RR),\CC)$. By passing to $(\mf g,K)$-modules, one receives an embedding
	\[D_k^+\into C^\infty((N(\RR),\psi_N)\backslash G(\RR))^{\mathrm{fin}}.\]
	To understand such maps, recall that a map out of $D_k^+$ is uniquely determined by where it sends the weight vector $v_k$, which satisfies $fv_k=0$. Thus, we are hunting for a smooth function on $G(\RR)$ which is left-invariant by $(N(\RR),\psi_N)$, right-invariant by $(\op{SO}(2),\chi_k)$, and killed by $f$. Well, the Iwasawa decomposition explains that
	\[N(\RR)\backslash\op{SL}_2(\RR)/\op{SO}_2(\RR)=\left\{\begin{bmatrix}
		a \\ & a^{-1}
	\end{bmatrix}:a>0\right\}.\]
	Thus, our smooth function can be described as a function on $\RR^+$. Lastly, we add in the condition that $fv_k=0$, which means that we are hunting for solutions to a first-order differential equation. Such a thing is unique (up to scalar), so there is a unique Whittaker functional (up to scalar).
\end{proof}
\begin{remark}
	For the discrete series, the Whittaker functional can be presented as given by functions on the upper-half plane, where it is given by the Fourier transform.
\end{remark}
\begin{remark}
	The same argument works for any irreducible $(\mf g,K)$-module of $\op{SL}_2(\RR)$, but in general, we cannot hunt for the condition $fv_k=0$. Indeed, for generic principal series representations, we are only characterized by the Casimir eigenvalue $\lambda$, so we are hunting for a function $F$ on $\RR^+$ satisfying $\Omega F=\lambda F$. This is now a second-order differential equation
\end{remark}
\begin{remark}
	It may seem confusing that $D_k^+$ has one Whittaker functional while the principal series has two. Roughly speaking, the problem is that the ``extra'' Whittaker functional for $D_k^+$ coming from the corresponding second-order differential equation $\Omega F=\lambda F$ does not actually always produce a copy of $D_k^+$. Indeed, there is a non-split extension of $D_k^+$ by the finite-dimensional representation with weight $k$ for which $fv_k\ne0$, and it has the same central Casimir eigenvalue.
\end{remark}
The moral is that one, in general, has two different Whittaker functionals. There are two ways to fix this problem. One way to fix this problem is to ask for moderate growth.
\begin{definition}[moderate growth] \label{def:mod-growth}
	Fix a split reductive group $G$ over $\RR$, and let $V$ be a continuous representation of $G$. Then a functional $\Lambda$ on $G$ has \textit{moderate growth} if and only if each $v\in V$ makes the function $G(\RR)\to\RR$ given by
	\[g\mapsto\Lambda(\pi(g)v)\]
	dominated by a polynomial in a choice of height of $g\in G(\RR)$.
\end{definition}
\begin{remark}
	By the Iwasawa decomposition, it is enough to restrict the moderate growth condition to the split torus of $G(\RR)$.
\end{remark}
\begin{proposition}
	Fix a split reductive group $G$ over $\RR$. Any $(\mf g,K)$-module $V$ admits at most one embedding into
	\[C^{\infty,\mathrm{moderate}}((N(\RR),\psi_N)\backslash G(\RR))^{\mathrm{fin}}.\]
\end{proposition}
\begin{proof}
	Omitted. This is a matter of solving some second-order differential equations.
\end{proof}
\begin{remark}
	Alternatively, one could ask for the functional $\Lambda\colon V\to\CC$ to be continuous. This is also unique (by a theorem of Shalika).
\end{remark}

\subsection{Global Whittaker Functionals}
Now that we have a finite and infinite story, we can tell a global story. Let's just straight from Whittaker functionals to Whittaker models.
\begin{definition}[Whittaker model]
	Fix a split reductive group $G$ over a global field $F$, and fix an automorphic representation $\pi$ of $G$. Let $\psi\colon F\backslash\AA_F\to S^1$ be a unitary character, and let $\psi_N\colon N(F)\backslash N(\AA_F)\to\CC$ to be the usual character given by summing the simple root spaces. Then a \textit{Whittaker model} is a $G(\AA_F)$-embedding
	\[\pi\into\mc A((N(\AA_F),\psi_N)\backslash G(\AA_F)),\]
	where the target is the set of $(N(\AA_F),\psi_N)$-invariant functions on $G(\AA_F)$, which are $Z\mf g_\infty$-finite, $K_\infty K_f$-finite, and have moderate growth.
\end{definition}
The uniqueness of the local Whittaker models implies that each $\pi$ has at most one Whittaker model, up to suitable scalars.
\begin{theorem} \label{thm:global-whittaker-unique}
	Fix a split reductive group $G$ over a global field $F$, and fix an automorphic representation $\pi$ of $G$. Then there is at most one subspace of $\mc A((N(\AA_F),\psi_N)\backslash G(\AA_F))$ isomorphic to $\pi$.
\end{theorem}
\begin{proof}
	We will work in the situation where $F$ is a function field so that we may avoid archimedean places. In this setting, we will actually be able to show that the Whittaker functional $\Lambda\colon V\to\CC$ is unique up to scalar. Well, at all but finitely many places in $S$, we see that $\pi_v$ is spherical and $\psi_v$ has conductor $\OO_v$ for $v\notin S$.

	Now, each $\pi_v$ for $v\notin S$ admits a unique Whittaker functional $\Lambda_v^\circ\colon V_v\to\CC$ up to scalar. Indeed, one can find an unramified character $\chi_v$ (in its Weyl orbit) for which $\pi_v$ is a subrepresentation of the principal series $\op{Ind}_B^G\chi_v$ for an unramified character $\chi_v$, which a unique nonzero Whittaker functional (up to scalar) by \Cref{prop:whittaker-principal-series}. We now normalize $\Lambda_v^\circ$ so that $\Lambda_v(\xi_v)=1$, where $\xi_v\in V_v$ is the distinguished spherical vector in the restricted tensor product $V=\bigotimes_v(V_v,\xi_v)$.

	We now return to $\Lambda$. Then $\Lambda(\eta)\ne0$ for some $\eta\in V$, and we may even assume that $\eta$ is a pure tensor because such elements span. Accordingly, we may write $\eta=\eta^{S'}\otimes\eta_{S'}$, where $\eta^{S'}=\bigotimes_{v\notin S'}\xi_v$ and $\eta_T\in\bigotimes_{v\in S'}V_v$; after possibly expanding $S$ and $S'$, we may assume $S=S'$. Thus, the restriction
	\[\Lambda\circ\left(\xi^S\otimes-\right)\colon\bigotimes_{s\in S}V_v\to\CC\]
	is a nonzero Whittaker functional for the $V_v$s. Now that we know they exist, we may go ahead and fix our own Whittaker functionals $\Lambda_v^\circ\colon V_v\to\CC$ as well, so we see that
	\[\Lambda\circ\left(\xi^S\otimes-\right)=c\bigotimes_{v\in S}\Lambda_v^\circ\]
	for some unique scalar $c$; notably, the finite tensor product of Whittaker representations continues to be multiplicity-free because it is the tensor product of such representations. Alternatively, one can proceed directly inductively: for just two places $v$ and $w$, we see $\Lambda(a_v,-)$ and $\Lambda^\circ(a_v,-)$ differ by a scalar for any fixed $a_v$, and this scalar is a Whittaker functional for the place $w$, so these scalars are all fixed. Then one can iterate this argument for finitely many places.

	We now claim that
	\[\Lambda\stackrel?=c\bigotimes_v\Lambda_v^\circ,\]
	which will complete the proof because the right-hand side does not depend on $\Lambda$. It is enough to check this equality on pure tensors, which may live in $\xi^T\otimes\bigotimes_{v\in T}V_v$ for a finite set $T$ containing $S$. Well, the two Whittaker functionals
	\[\Lambda\circ\left(\xi^T\otimes-\right)\qquad\text{and}\qquad c\Lambda^\circ\left(\xi^T\otimes-\right)\]
	on $\bigotimes_{v\in T}V_v$ must agree up to a scalar, by the above argument, and this scalar is fixed by computing both sides on $\eta$!
\end{proof}

\end{document}